\section{从 Blocklist 到 Neurosymbolic Hybrid}

最后一组 slides 不只是列 follow-up papers，而是在回答“一个 imperfect judgment model 能否作为系统组件”。应用价值取决于把 Delphi 放在哪里、它能否被 override、怎样表达 identity 与 cultural variation，以及 top-down principle 如何约束 bottom-up predictions。

\subsection{Keyword Blocklist 为什么不够}

Alexa Prize slide 先庆祝冠军，随后回顾 2017 年的安全做法：用一组 sensitive keywords 阻止系统讨论性、同性恋、种族等话题。blocklist 易实现，却把复杂 context 简化为词表；它会同时漏掉不含关键词的伤害表达，也会误伤正常、求助或教育性讨论。

\lecturefigure{slide-45-alexa-prize-sensitive-keywords.jpg}{Alexa Prize 的 sensitive-keyword blocklist 与其局限}{官方录像恢复幻灯片 V045；Stanford CS25 Lecture 16}

\teachervoice{讲者把 blocklist 称为当年的 status quo，并问 2021 年是否还只能这样做。她没有主张 Delphi 可以直接替代 safety policy，而是主张系统至少需要理解“相同词汇在不同情境中为何不同”。}

\subsection{Identity-Specific Judgment：Context 增加，风险也增加}

identity slide 比较同一关系或行为在不同 gender / sexuality 描述下的判断。引入 identity 可以纠正“一刀切”模型，也可能放大 stereotype。读图时应检查 model 是否真正使用 relevant context，还是把 identity token 当作 shortcut；评价需要 counterfactual pairs 与 group-level error analysis。

\lecturefigure{slide-46-identity-specific-moral-judgments.jpg}{Delphi 是否能做 identity-specific moral judgment}{官方录像恢复幻灯片 V046；Stanford CS25 Lecture 16}

\subsection{Controllable Generation 与 Social Norm Alignment}

follow-up slide 左侧展示 controllable text generation：给定 persona、norm 或 constraint，引导系统生成更合适的故事或对话；右侧文字明确指出现有 dialogue systems 对冒犯内容可能盲目同意，目标是让 agent 学习 human moral preference。功能越强，越需要报告 preference 来自谁。

\lecturefigure{slide-47-controllable-generation-bias.jpg}{用 Delphi 作为 controllable generation 的 commonsense / norm component}{官方录像恢复幻灯片 V047；Stanford CS25 Lecture 16}

下一页展示 social norms 与 values 在 interactive narratives 中的对齐，并引用游戏 agent 研究。图中的流程把 narrative state、value model 与 action choice连接起来。它支持“judgment model 可作为 reward / filter”这一系统角色，却不能保证 long-horizon agent 不会通过 proxy optimization 绕过约束。

\lecturefigure{slide-48-align-social-norms-interactive-narratives.jpg}{在互动叙事中对齐 social norms 与 values}{官方录像恢复幻灯片 V048；Stanford CS25 Lecture 16}

\subsection{研究议程：Ethically、Socially、Culturally Aware}

路线图把 ethically-informed、socially-aware、culturally-inclusive 三个要求放在同一条道路上，并把 interpretability 与 diverse viewpoints 的 dynamic customization 作为 next priority。三个词不是同义反复：伦理关注 harm 与 principle，社会关注关系与制度，文化关注分布差异和谁被代表。

\lecturefigure{slide-49-ethical-social-cultural-ai.jpg}{Machine ethics 的下一阶段：可解释、动态、文化多元}{官方录像恢复幻灯片 V049；Stanford CS25 Lecture 16}

\subsection{Delphi Hybrid：为什么要回到 Neurosymbolic}

hybrid title slide 宣布“Demonstrative Neuro-symbolic Hybrid Moral Reasoning System”。回到整堂课的结构，neural component 提供 language 与 commonsense generation，symbolic / principle component 提供 top-down constraints。它不是把规则盖在模型输出上，而是尝试让案例直觉和原则互相校验。

\lecturefigure{slide-50-hybrid-moral-reasoning.jpg}{Delphi Hybrid：示范性的 neuro-symbolic moral reasoning system}{官方录像恢复幻灯片 V050；Stanford CS25 Lecture 16}

原始 Delphi failures slide 写出两类缺陷：lack of commonsense knowledge 与 lack of interpretability。一个例子中模型对“genocide if it creates jobs”给出错误判断，表明 surface association 会压过原则；缺乏 explanation 又使错误难以定位。hybrid 的目标是同时增加 knowledge path 与 constraint path。

\lecturefigure{slide-51-original-delphi-failures.jpg}{原始 Delphi 的常识缺口与不可解释性}{官方录像恢复幻灯片 V051；Stanford CS25 Lecture 16}

\subsection{Bottom-Up 与 Top-Down 的理论框架}

理论 slide 引用 John Rawls 的 reflective equilibrium：bottom-up human judgments 提供案例直觉，top-down constraints 提供原则，两端需要反复修正。Delphi 原系统主要位于 bottom-up 一侧；hybrid 要问的是如何把 moral principle 编码成可参与 inference 的 symbolic constraints。

\lecturefigure{slide-52-theoretical-framework.jpg}{Reflective equilibrium：bottom-up intuition 与 top-down constraints}{官方录像恢复幻灯片 V052；Stanford CS25 Lecture 16}

下一页用 Bernard Gert 的 common morality 与十条 moral rules 作为 top-down 候选，例如不要杀人、不要欺骗、遵守承诺等。规则表看似清晰，却仍需要 context、exception 与 conflict resolution。图的教学价值在于提供可显式审计的 principle vocabulary，而不是宣称十条规则能机械解决所有伦理问题。

\lecturefigure{slide-53-top-down-common-morality.jpg}{Top-down moral principles：common morality 与十条规则}{官方录像恢复幻灯片 V053；Stanford CS25 Lecture 16}

\subsection{Hybrid Architecture 与约束求解}

详细架构 slide 把 situation 输入、rule retrieval、commonsense generation、Maieutic graph 与 final judgment 串起来。读图时应沿数据流从左到右：先生成 candidate reasons，再检索或激活原则，随后用图结构检查 contradiction，最终输出 judgment 与 rationale。任一模块错误都会传播，因此系统需要 module-level evaluation。

\lecturefigure{slide-54-hybrid-architecture-detail.jpg}{Delphi Hybrid 的 retrieval、generation、graph 与 judgment pipeline}{官方录像恢复幻灯片 V054；Stanford CS25 Lecture 16}

可以用一个简化目标表示 hybrid inference：
\begin{equation}
z^*=\arg\max_z
\left[
\sum_i \alpha_i s_i^{\mathrm{neural}}(z)
+\sum_j \beta_j\mathbf{1}[C_j(z)]
\right],
\end{equation}
其中 $s_i^{\mathrm{neural}}$ 是 commonsense / norm model 的证据分数，$C_j$ 是 top-down moral constraints，$\alpha_i$ 与 $\beta_j$ 控制两类证据权重。公式只是教学抽象；真实系统还要处理规则冲突、检索失败和 uncertain abstention。

\begin{lstlisting}[caption={Neuro-symbolic moral reasoning 的模块化循环}]
def hybrid_judgment(situation):
    candidate_reasons = commonsense_model.generate(situation)
    principles = retrieve_moral_rules(situation, candidate_reasons)
    graph = build_argument_graph(situation, candidate_reasons, principles)
    judgment, rationale = constrained_inference(graph)
    if low_confidence(judgment) or unresolved_conflict(graph):
        return abstain_with_audit_trace(graph)
    return judgment, rationale
\end{lstlisting}

\subsection{Hybrid 的结果与限制}

“Where does the original Delphi fail?”页把 Maieutic graph 与 architecture 缺口标回原系统，说明改进点对应已知 failure，而不是随意增加模块。它要求 evaluator 分别测 commonsense retrieval、principle coverage、graph consistency 与 final judgment，避免一个总分掩盖 pipeline bottleneck。

\lecturefigure{slide-55-where-delphi-fails.jpg}{把 hybrid 模块映射回原始 Delphi failure modes}{官方录像恢复幻灯片 V055；Stanford CS25 Lecture 16}

最后的 adversarial result bars 比较 Delphi 与 Delphi-hybrid 在 overall、certain、ambiguous queries 上的表现。hybrid 在课堂展示中更高，旁边注明将 keywords 与 interpretable hybrid moral reasoning 配对。读图时要注意 test set 仍有限，提升不能证明 top-down rules 已覆盖开放世界，也不能忽略额外 inference latency。

\lecturefigure{slide-56-adversarial-results-hybrid.jpg}{Delphi-hybrid 在 adversarial open-ended moral queries 上的结果}{官方录像恢复幻灯片 V056；Stanford CS25 Lecture 16}

\begin{warningbox}{Hybrid 不是“神经 + 符号 = 自动安全”}
模块化提高可检查性，也增加错误接口。retriever 可能选错规则，generator 可能编造理由，solver 可能在错误 clauses 上得到一致答案。安全收益来自可观测、可分解、可拒答和可人工复核，而不是 architecture 标签本身。
\end{warningbox}

\subsection{本章小结}

从 keyword blocklist 到 hybrid，系统逐步增加 context、identity、norm 与 principle。更强的伦理组件必须同时增加审计能力：记录数据来源、显示 rationale provenance、允许 abstention、支持 override，并把 authority 留在制度与人类责任链中。

